\section{Current repair, uncertainty, and target-error diagnostics}\label{ec:r12-diagnostics}
This section gives the complete diagnostics supporting Section~\ref{sec:r12-experiment}. The frozen design uses four training seeds (12000--12003), twelve test seeds (920000--920011) in each of four graph families, and no selection based on a favorable seed. The independent-instance sample size is 48, not the number of seed--instance--method combinations. The graph families are fixed strata. Bootstrap intervals resample seeds jointly across instances and instances within strata; they do not resample individual actor outcomes as independent observations. The complete original ten-seed/eight-instance study remains above and is not pooled with this one.

\subsection{Pre-repair violations, projection, and feasible actors}
Table~\ref{tab:r12-repair} distinguishes the actor before continuation repair from each repaired policy. All gains use the same objective. A common radial factor can modify every coordinate because one continuation row is violated; its cost can therefore extend beyond the node causing the violation. Weighted projection is solved by SLSQP with analytic derivatives of squared distance, followed by rational rounding and exact subtree repair. Any residual rounding correction is charged to the projection pipeline. The raw records retain both projection solver status and all exact feasibility checks. Neither a projection status nor a neural output is used as a certificate.

The tree-transfer actor fits nonnegative flows to the predicted displacement by weighted nonnegative least squares. For a transfer on edge $u\to v$ and service $j$, the direction has components $1/(w_ua_{uj})$ at the parent and $-1/(w_va_{vj})$ at the child. A positive combination has subtree payment change equal to minus the entering transfer, so all continuation caps remain satisfied. The final scalar step obeys every tier boundary. At a boundary, this step can be zero; the implementation records that outcome. A simpler segment from the comparator toward the clipped actor is also recorded, but can be identically static when every cap initially binds. Neither parameterization is claimed to span all vector-service accepted improvements, because within-node cross-service reallocations may also matter.
\begin{table}[p]\centering
\caption{Raw actors and alternative acceptance repairs}\label{tab:r12-repair}
\setlength{\tabcolsep}{3.5pt}\begin{tabular}{llrrrrrr}
\toprule
Family & Actor & \shortstack{Feasible\\before} & \shortstack{Mean\\scale} & \shortstack{Before\\gain} & \shortstack{Radial\\gain} & \shortstack{Projection\\gain} & \shortstack{Tree-flow\\gain} \\
\midrule
Cycle & Value & 9/48 & 0.933 & -0.8182 & -0.8228 & -0.8182 & 0.0023 \\
Cycle & Value--gradient & 13/48 & 0.953 & -0.7616 & -0.7664 & -0.7620 & 0.0023 \\
Random-edge & Value & 8/48 & 0.937 & -0.8192 & -0.8215 & -0.8189 & 0.0029 \\
Random-edge & Value--gradient & 8/48 & 0.959 & -0.7621 & -0.7642 & -0.7622 & 0.0029 \\
Geometric & Value & 9/48 & 0.950 & -1.4773 & -1.4803 & -1.4765 & 0.0015 \\
Geometric & Value--gradient & 10/48 & 0.971 & -1.4383 & -1.4407 & -1.4383 & 0.0018 \\
Weighted block & Value & 8/48 & 0.956 & -2.0912 & -2.0952 & -2.0903 & -0.0036 \\
Weighted block & Value--gradient & 9/48 & 0.976 & -2.0669 & -2.0703 & -2.0668 & -0.0023 \\
\bottomrule
\end{tabular}
\begin{minipage}{.98\textwidth}\small
Before denotes box clipping without continuation repair. Negative raw outcomes are retained. Projection minimizes weighted Euclidean actor distance, not economic loss. Tree-flow uses nonnegative child-to-parent payment transfers and a box-limited step. All reported repaired policies receive the same exact feasibility and value audit. All gains are relative to the reoptimized static protocol per discounted service--review; the comparison is not field-calibrated.
\end{minipage}
\end{table}


\subsection{Gate ratios and measured large-instance fallback}
The denominator in Table~\ref{tab:r12-ratios} is the best recorded classical feasible gain on the corresponding original instance, not an unknown exact optimum. Its tiny certified gap is retained in the predecessor results. The proposal gap uses the updated residual calculation; the selected gap also uses a separately audited static upper bound. Both are valid upper bounds on the selected policy's loss. The original weighted-block proposals all select static before refinement. The current fallback runs the cached Newton method afresh and audits its full contingent policy. It does not reuse a fictional learned Hessian or count the fallback as a neural success. The full timing record separates this measurement from the original cached-quadratic measurements.
\begin{table}[p]\centering
\caption{Original nonlinear transfer: certificate size relative to available gain}\label{tab:r12-ratios}
\setlength{\tabcolsep}{3.5pt}\begin{tabular}{llrrr}
\toprule
Family & Actor & \shortstack{Mean raw\\gain} & \shortstack{Proposal gap /\\classical gain} & \shortstack{Selected gap /\\classical gain} \\
\midrule
Cycle & Value & 0.008528 & 0.59 & 0.59 \\
Cycle & Value--gradient & 0.009686 & 0.44 & 0.44 \\
Random-edge & Value & 0.012158 & 0.43 & 0.43 \\
Random-edge & Value--gradient & 0.013160 & 0.24 & 0.24 \\
Geometric & Value & 0.002601 & 3.00 & 3.00 \\
Geometric & Value--gradient & 0.003773 & 2.31 & 2.31 \\
Weighted block & Value & -0.097068 & 1340.86 & 23.07 \\
Weighted block & Value--gradient & -0.098095 & 1484.75 & 23.07 \\
\bottomrule
\end{tabular}
\begin{minipage}{.98\textwidth}\small
Ratios are medians across the original two instances and ten training seeds, using the classical certified feasible gain as denominator. The selected gap uses both the proposal and static upper bounds, as in the main-paper gate equation. The weighted-block gate still has a large gap; a no-harm choice is not an accuracy certificate. All forty weighted-block pipelines were subsequently refined by the cached classical solver and independently audited.
\end{minipage}
\end{table}


\subsection{Timing, common setup, and endpoint errors}
Table~\ref{tab:r12-repeated} uses five repetitions after one warmup on a separately identified first instance per family, with pipeline order randomized within repetition. The common outside-contract solve, sparse operator setup, and cached static certificate are listed in the raw setup file. The table includes all online construction and refinement stages. Python, NumPy, SciPy, operating-system, and thread settings are recorded with the results. The numerical stopping flag is not substituted for the common rational error threshold.
\begin{table}[p]\centering
\caption{Timing variation separated from test-instance and training-seed variation}\label{tab:r12-repeated}
\setlength{\tabcolsep}{3.5pt}\begin{tabular}{llrrrr}
\toprule
Family & Pipeline & Repetitions & Minimum & Median & Maximum \\
\midrule
Cycle & Cold SLSQP & 5 & 11.73 & 11.76 & 11.85 \\
Cycle & Value & 5 & 18.92 & 19.10 & 19.50 \\
Cycle & Value--gradient & 5 & 18.53 & 18.74 & 18.95 \\
Random-edge & Cold SLSQP & 5 & 41.13 & 41.73 & 42.01 \\
Random-edge & Value & 5 & 48.70 & 49.65 & 59.45 \\
Random-edge & Value--gradient & 5 & 47.67 & 48.55 & 52.80 \\
Geometric & Cold SLSQP & 5 & 67.37 & 67.63 & 68.11 \\
Geometric & Value & 5 & 91.91 & 93.55 & 93.83 \\
Geometric & Value--gradient & 5 & 91.43 & 92.89 & 93.77 \\
Weighted block & Cold SLSQP & 5 & 108.66 & 109.76 & 110.65 \\
Weighted block & Value & 5 & 138.82 & 139.56 & 140.86 \\
Weighted block & Value--gradient & 5 & 135.83 & 137.49 & 138.97 \\
\bottomrule
\end{tabular}
\begin{minipage}{.98\textwidth}\small
Milliseconds to the same $10^{-7}$ certified tolerance on the first new instance of each family, one frozen training seed, five repetitions after one warmup. Pipeline order is randomized within repetition. This separate experiment does not enlarge the economic test sample; every timing record and final policy is retained.
\end{minipage}
\end{table}


The endpoint sensitivity study regenerates twelve independent training-sized problems and solves two endpoint value queries at each of three step sizes. Tariffs, outside protocol, and participation caps are fixed during each pair. The reported finite-difference error bound is the bias term from the forcing-value gradient Lipschitz constant plus the sum of endpoint certified gaps divided by twice the step size. Additional endpoint-error allowances in Table~\ref{tab:r12-fd} are an analytical sensitivity calculation, explicitly distinguished from newly simulated label noise. The smaller step reduces the bias bound while amplifying oracle uncertainty. A small observed mean improvement between training losses is not itself evidence that the directional targets are sufficiently accurate.
\begin{table}[p]\centering
\caption{Directional targets: step size and certified oracle error}\label{tab:r12-fd}
\setlength{\tabcolsep}{3.5pt}\begin{tabular}{rrrr}
\toprule
Step size & \shortstack{Additional endpoint\\error allowance} & \shortstack{Maximum measured\\direction error} & \shortstack{Maximum certified\\target-error bound} \\
\midrule
$1/80$ & 0e+00 & 5.58e-07 & 0.00208 \\
$1/80$ & 1e-07 & 5.58e-07 & 0.00209 \\
$1/80$ & 1e-05 & 5.58e-07 & 0.00288 \\
$1/40$ & 0e+00 & 8.89e-06 & 0.00417 \\
$1/40$ & 1e-07 & 8.89e-06 & 0.00417 \\
$1/40$ & 1e-05 & 8.89e-06 & 0.00457 \\
$1/20$ & 0e+00 & 3.45e-05 & 0.00833 \\
$1/20$ & 1e-07 & 3.45e-05 & 0.00834 \\
$1/20$ & 1e-05 & 3.45e-05 & 0.00853 \\
\bottomrule
\end{tabular}
\begin{minipage}{.98\textwidth}\small
Twelve independently generated training-sized problems. Endpoint values are re-solved and rationally certified at every step size. The extra error column is an explicit sensitivity allowance added to each endpoint certificate, not a claim that fresh noisy labels were simulated. Measured error compares with the separately solved central actor and is numerical; the displayed bound includes oracle gaps and the finite-difference bias bound.
\end{minipage}
\end{table}


\subsection{Structural cross-checks and their status}
The proof of the capacity theorem does not depend on numerical testing. Independent LP formulations nevertheless cross-check its directional and normal-cone forms on 180 randomly generated problems with boundary coordinates, zero and signed commitment coefficients, several constraints, and nonconstant comparators. A separate set of 48 trees, with 3, 7, 15, and 31 nodes, compares the full balance program with the cumulative-cut critical-friction program. Thresholds are also tested just below, at, and above the returned value. These are floating-point HiGHS diagnostics, not exact theorem certificates. The two-review threshold $72/155$, the nonconstant boundary optimality interval $[0,1]$, and the first-order boundary-loss example are independently evaluated with rational arithmetic. Initial installation and outgoing switching edges are included throughout.
